\chapter{Storage, Transport and Physical Optionality}\label{ch:s2:storage-transport-and-physical-optionality}

A trader who leases a gas storage cavern owns a string of calendar spread options; one who holds pipeline capacity owns \omterm{def:s2:commodity-spreads:location}{location spread} options. The
physical asset is valued as the options it contains, and the trader's job is to turn them into money with hedges. Secomandi found that most of the
value of adapting to price uncertainty can be generated simply by re-optimising, again and again, a plan that ignores the uncertainty. On this
chapter's synthetic lease, locking today's best plan is worth \$716,000. Re-optimising it each month on the new forward curve and re-hedging adds
\$56,800 on average, and on no path does it earn less. At two cents per unit of forward traded it adds \$26,600 and loses to the locked plan on a third
of paths. The build is \texttt{firm.physopt}.

\section{Storage as an option}

\begin{definition}[Storage deal]\label{def:s2:storage-transport-and-physical-optionality:storage}
A \emph{storage deal}\index{storage deal} is a lease of storage capacity (working volume and injection and withdrawal rates, for a period) that the
trader monetises by buying the commodity when it is cheap, holding it and selling when it is dear, hedging each planned injection and withdrawal in the
forward market.
\end{definition}

The lease in \texttt{firm.physopt} holds 10 units of 100,000 MMBtu, injects up to 2 units and withdraws up to 4 a month, and costs 2 cents per MMBtu
each way. The forward curve and the price model are Book~6's (chapter~16): an illustrative Henry Hub-like curve from April (\$2.60) to January
(\$3.50) and a two-factor model with 60\% short-end and 20\% long-end volatility. The intrinsic plan buys 2 units in each month from April to August
and sells 4 in December, 4 in January and 2 in February. Locked with forwards, it is worth \$716,000.

Four valuation methods give four numbers:

\begin{center}
\small
\begin{tabular}{@{}lr@{}}
\toprule
method & value (\$ thousands)\\
\midrule
intrinsic: today's best plan, locked & 716\\
spread options: the plan's ten injection-withdrawal pairs, each an option & 722\\
least-squares Monte Carlo on spot (Book~6) & 760\\
rolling intrinsic: re-optimised each month, hedges traded & 773\\
\bottomrule
\end{tabular}
\end{center}

The spread-option method pairs each unit injected with a unit withdrawn, first in first out, and values each pair as a calendar spread option on the
two forwards struck at the operating costs (\cref{lst:s2:storage-transport-and-physical-optionality:options}). Here the pairs are deep in the money
and add little. Least-squares Monte Carlo acts on spot prices alone and estimates continuation values by regression, which loses a little.
Rolling intrinsic uses the whole forward curve each month. Lai, Margot and Secomandi found that practitioners' heuristics are fast but significantly
suboptimal, and that with periodic re-optimisation inside a Monte Carlo simulation they become nearly optimal, with one exception.

\section{Monetising the options}

The value is realised through a hedge programme (\cref{lst:s2:storage-transport-and-physical-optionality:programme}). The static programme buys
forwards for every planned injection and sells forwards for every planned withdrawal, 20 units in all, and delivers the plan at spot. Whatever the
prices do, it earns \$716,000: the hedges' gains and losses exactly offset the spot prices paid and received. The rolling programme re-optimises the
plan each month on the simulated forward curve, given the gas already in store, and trades the difference in hedges. A new plan is chosen only if
it is worth more on the current curve than the old one, and the hedges lock that gain; so, without trading costs, rolling can never earn less than
the static programme on any path.

\begin{center}
\small
\begin{tabular}{@{}lcccc@{}}
\toprule
cost per unit of forward traded & 0 & 0.5 cent & 2 cents & 2 cents, gated\\
\midrule
extra over the static programme (\$ thousands) & 56.8 & 49.2 & 26.6 & 37.4\\
paths below the static programme & 0\% & 3.1\% & 32.6\% & 0\%\\
extra units of forward traded & 15.1 & 15.1 & 15.1 & 9.2\\
\bottomrule
\end{tabular}
\end{center}

The static programme also pays for its own 20 units of hedges, so the comparison is at equal costs. At two cents per unit the rolling programme
still earns more on average, but it loses to the static one on a third of paths, because it re-hedges whenever the new plan is better by any amount
(\cref{fig:s2:storage-transport-and-physical-optionality:rolling}). The gated version re-hedges only when the gain on the curve exceeds the cost of the
trades. It trades 9.2 extra units instead of 15.1, earns \$37,400 more than the static programme on average and never less.

\begin{omfigure}
\begin{tikzpicture}
\begin{axis}[width=11cm,height=5.4cm,xlabel={\small rolling minus static programme (\$ per MMBtu-unit; \$100,000 per unit)},
  ylabel={\small share of paths (\%)},tick label style={font=\footnotesize},xmin=-1,xmax=3,ymin=0,ymax=26,grid=major,grid style={gray!25},
  table/col sep=comma,legend style={font=\scriptsize,at={(0.97,0.97)},anchor=north east}]
\addplot[const plot,omThm,thick] table[x=lo,y=free]{figdata/strategies-2/22-storage-transport-and-physical-optionality/rolling.csv};
\addplot[const plot,omDef,thick,dashed] table[x=lo,y=cost2]{figdata/strategies-2/22-storage-transport-and-physical-optionality/rolling.csv};
\legend{no trading cost,2 cents per unit}
\end{axis}
\end{tikzpicture}

\omcaption{Distribution over 2,000 simulated paths of the rolling-intrinsic programme's P\&L less the static programme's, without trading costs and
at 2 cents per unit of forward traded; the end bins collect the tails. Data: \texttt{s2\_physopt.programmes}.}
\label{fig:s2:storage-transport-and-physical-optionality:rolling}
\end{omfigure}

\section{Transport as an option}

\begin{definition}[Transport option]\label{def:s2:storage-transport-and-physical-optionality:transport}
A \emph{transport option}\index{transport option} is the right, given by pipeline or shipping capacity, to move a commodity from one location to
another in a period at a variable cost; it is worth a spread option on the destination price less the origin price, struck at that cost.
\end{definition}

Capacity of 1 unit a month from the hub to a market-area hub, whose forwards stand above the hub's by 5 cents in spring to 60 cents in January, with a
variable cost of 10 cents, is worth \$139,000 on today's forwards: only the winter months pay. As a strip of spread options, with volatilities of 45\%
and 55\% and a correlation of 0.85, it is worth \$369,000 (\cref{fig:s2:storage-transport-and-physical-optionality:transport}). Its value depends on
how independently the two hubs move: \$580,000 at a correlation of 0.5 and \$273,000 at 0.95. The holder monetises it like storage: by selling the
spread forward in the months it is in the money, and re-optimising as the basis moves.

\begin{omfigure}
\begin{tikzpicture}
\begin{axis}[width=11cm,height=5.2cm,ybar,bar width=6pt,xlabel={\small delivery month},ylabel={\small \$ per MMBtu},
  tick label style={font=\footnotesize},ymin=0,ymax=0.8,xtick=data,xticklabels={Apr,May,Jun,Jul,Aug,Sep,Oct,Nov,Dec,Jan,Feb,Mar},
  grid=major,grid style={gray!25},table/col sep=comma,legend style={font=\scriptsize,at={(0.03,0.97)},anchor=north west}]
\addplot[fill=gray!50,draw=gray] table[x=k,y=intrinsic]{figdata/strategies-2/22-storage-transport-and-physical-optionality/transport.csv};
\addplot[fill=omThm!60,draw=omThm] table[x=k,y=option]{figdata/strategies-2/22-storage-transport-and-physical-optionality/transport.csv};
\legend{intrinsic,spread option}
\end{axis}
\end{tikzpicture}

\omcaption{A year of monthly transport capacity from the hub to a market-area hub: the intrinsic value on today's forwards and the spread-option value
of each month, per MMBtu. Data: \texttt{s2\_physopt.transport\_strip}.}\label{fig:s2:storage-transport-and-physical-optionality:transport}
\end{omfigure}

A cargo of liquefied natural gas committed to Europe at \$10.00 per MMBtu that can instead go to Asia at \$10.50, for 80 cents more in shipping,
has no intrinsic value in switching. The switch is still an option on the Asian price less the European price, struck at 80 cents: with three months
to decide, volatilities of 60\% and a correlation of 0.8, it is worth 64 cents per MMBtu, and \$1.08 at a correlation of 0.5.

\section{What the asset-backed trader owns}

\begin{definition}[Asset-backed trading]\label{def:s2:storage-transport-and-physical-optionality:assetbacked}
\emph{Asset-backed trading}\index{asset-backed trading} is trading around physical assets that the firm owns or leases (storage, transport capacity,
generation, cargoes), valuing them as options and monetising their flexibility with hedges in the forward and options markets.
\end{definition}

The asset-backed trader owns three things a financial trader does not. The first is the options themselves, which cannot be bought on an exchange.
The second is the obligation to operate them, with the physical constraints that make the options valuable and the plans complicated. The third is
information: flows through the asset show supply and demand before prices do. What the trader does not own is certainty that the spreads will be there.
From 2007 to 2025, the Henry Hub winter spot price (December to February) was above the preceding summer's (April to June) in 12 of 19 years. It
averaged 7 cents less, with a standard deviation of \$2.04. In 2008 the winter was \$6.12 below the summer, and in 2022 \$3.71. An operator who injected
and waited without hedging earned those numbers. The hedge programme exists to replace them with the forward spread at the time of the decision.

\section{Strategy files}

\begin{strategyfile}{Rolling-intrinsic storage}\label{strat:s2:storage-transport-and-physical-optionality:rolling}
\sfield{Who pays you, and why} Consumers who pay a winter premium; the forward curve's seasonal spread pays for holding gas through the summer.
\sfield{Instruments and venues} A storage lease; monthly gas forwards or futures.
\sfield{Signal} The best plan on the current forward curve.
\sfield{Sizing and execution} Hedge the whole plan; re-optimise monthly or on large moves; re-hedge only when the gain exceeds the cost.
\sfield{Costs} The lease; injection and withdrawal fees; forward bid-ask.
\sfield{How it dies} Lease costs above the value; flat curves; operational outages.
\sfield{Horizon, capacity, infrastructure} A storage year; scheduling and nominations.
\sfield{Backtest honestly} Forward curves as they were; physical constraints enforced; costs on every re-hedge.
\sfield{Sources} Secomandi (2010); Lai, Margot and Secomandi (2010); this chapter: \$56,800 above intrinsic without costs.
\end{strategyfile}

\begin{strategyfile}{Spread-option storage hedge}\label{strat:s2:storage-transport-and-physical-optionality:options}
\sfield{Who pays you, and why} Buyers of calendar spread options who want the flexibility the storage owner has.
\sfield{Instruments and venues} Calendar spread options on gas futures; over-the-counter spread options.
\sfield{Signal} The storage's option value against the options' market price.
\sfield{Sizing and execution} Sell the options the facility can deliver: pairs of injection and withdrawal months.
\sfield{Costs} Option bid-ask.
\sfield{How it dies} Selling more optionality than the facility's rates allow.
\sfield{Horizon, capacity, infrastructure} A season.
\sfield{Backtest honestly} Exercise constrained by the facility's rates.
\sfield{Sources} This chapter: the pairs are worth \$722,000 against \$716,000 intrinsic.
\end{strategyfile}

\begin{strategyfile}{Transport capacity optimisation}\label{strat:s2:storage-transport-and-physical-optionality:transport}
\sfield{Who pays you, and why} Buyers in constrained market areas; the basis pays for the pipe.
\sfield{Instruments and venues} Pipeline capacity; basis swaps or futures at both hubs.
\sfield{Signal} The forward basis against the variable cost.
\sfield{Sizing and execution} Sell the basis forward in months in the money; re-optimise as it moves.
\sfield{Costs} Capacity charges; fuel; basis bid-ask.
\sfield{How it dies} New pipelines that close the basis; correlation that rises.
\sfield{Horizon, capacity, infrastructure} Seasons to years; nominations.
\sfield{Backtest honestly} Basis at both hubs as it was; capacity outages.
\sfield{Sources} This chapter: \$139,000 intrinsic, \$369,000 as options.
\end{strategyfile}

\begin{strategyfile}{Cargo diversion}\label{strat:s2:storage-transport-and-physical-optionality:cargo}
\sfield{Who pays you, and why} Markets that bid for flexible supply when tight.
\sfield{Instruments and venues} LNG cargoes with destination flexibility; regional gas futures; freight.
\sfield{Signal} Netbacks by destination: price less shipping and regasification.
\sfield{Sizing and execution} Hedge the committed destination; hold the switch as an option.
\sfield{Costs} Shipping, boil-off, canal and port charges.
\sfield{How it dies} Contract restrictions; freight spikes; converging regional prices.
\sfield{Horizon, capacity, infrastructure} Weeks to months; shipping and scheduling.
\sfield{Backtest honestly} Freight and netbacks at the decision date.
\sfield{Sources} This chapter: 64 cents per MMBtu at a correlation of 0.8.
\end{strategyfile}

\section{Tutorial: the asset as options}\label{tut:s2:storage-transport-and-physical-optionality:tut}

\begin{tutorial}
\textbf{Goal.} Value a storage lease four ways, run its static and rolling hedge programmes path by path with trading costs, and value transport
capacity and a cargo switch as spread options. \textbf{End state:} the two tables and two figures.
\begin{tutsteps}
\item \textbf{The hedge programme}.
\omcode{code/firm/physopt/firm_physopt.py}{60}{91}{Locking the plan, re-optimising, and re-hedging when it pays.}\label{lst:s2:storage-transport-and-physical-optionality:programme}
\item \textbf{The options}.
\omcode{code/firm/physopt/firm_physopt.py}{94}{124}{Spread-option pairs, transport and diversion.}\label{lst:s2:storage-transport-and-physical-optionality:options}
\item \textbf{Run} \texttt{storage()}, \texttt{hedging()}, \texttt{gated()}, \texttt{transport\_strip()}, \texttt{cargo()},
\texttt{fig\_physopt.py} and, once, \texttt{s2\_fetch\_hh.py}.
\end{tutsteps}
\textbf{What to change next.} Re-optimise only on large curve moves; add a lease price and find the break-even; hedge the transport strip and
measure its P\&L path by path.
\end{tutorial}

\section{Build: physical optionality}\label{bld:s2:storage-transport-and-physical-optionality:physopt}

\begin{build}
\bfield{Purpose} Storage hedge programmes path by path, spread-option decomposition, transport strips and cargo switches, on Book 6's energy model.
\bfield{Interface} \texttt{plans(fac, prices, inv)}, \texttt{hedge\_programme(fac, model, F0, T, paths, seed, cost, roll, gate)},
\texttt{spread\_option\_pairs(fac, F0, T, model)}, \texttt{transport(FA, FB, T, cost, sA, sB, rho)}, \texttt{diversion(F1, F2, K, T, s1, s2,
rho)}.
\bfield{Rules} Hedges long for injections and short for withdrawals; marks at each month's curve; physical at spot; costs per unit of forward traded.
\bfield{Acceptance tests} \texttt{code/firm/physopt/tests/}: plans equal Book 6's intrinsic; the static programme earns the intrinsic value on every
path and the rolling one never less without costs; pairs, transport and diversion by hand.
\bfield{Stretch} Ratchets (rates depending on inventory); lease pricing; swing contracts.
\end{build}

\begin{omsources}
\item N.~Secomandi, ``Optimal commodity trading with a capacitated storage asset'', \emph{Management Science} 56(3), 2010.
\item G.~Lai, F.~Margot and N.~Secomandi, ``An approximate dynamic programming approach to benchmark practice-based heuristics for natural gas
storage valuation'', \emph{Operations Research} 58(3), 2010.
\item US Energy Information Administration, Henry Hub natural gas spot price, daily, via FRED, accessed 25 September 2026.
\end{omsources}

\section{Exercises}

\begin{exercise}[$\star$]\label{exo:s2:storage-transport-and-physical-optionality:1}
Compute the intrinsic plan's value: 2 units bought each month from April to August at \$2.60, \$2.62, \$2.68, \$2.74 and \$2.78; 4 sold at \$3.45 and
4 at \$3.50; 2 at \$3.30; 2 cents per MMBtu each way; units of 100,000 MMBtu.
\end{exercise}

\begin{exercise}[$\star$]\label{exo:s2:storage-transport-and-physical-optionality:2}
In January the market-area hub is 60 cents above the hub and transport costs 10 cents. What is one unit of January capacity worth on the forwards?
\end{exercise}

\begin{exercise}[$\star$]\label{exo:s2:storage-transport-and-physical-optionality:3}
A cargo committed to Europe at \$10.00 can go to Asia for 80 cents more in shipping. What is the switch's intrinsic value when Asia is at \$10.50,
and at \$11.20?
\end{exercise}

\begin{exercise}[$\star\star$]\label{exo:s2:storage-transport-and-physical-optionality:4}
Why can the rolling programme never earn less than the static one on any path when trading is free?
\end{exercise}

\begin{exercise}[$\star\star$]\label{exo:s2:storage-transport-and-physical-optionality:5}
Why does re-optimising a plan that ignores uncertainty capture most of the option value?
\end{exercise}

\begin{exercise}[$\star\star$]\label{exo:s2:storage-transport-and-physical-optionality:6}
Why is transport capacity worth more when the two hubs are less correlated?
\end{exercise}

\begin{exercise}[$\star\star\star$]\label{exo:s2:storage-transport-and-physical-optionality:7}
\emph{Coding.} Run \texttt{gated(0.005)}. How does gating change the rolling programme at half a cent per unit?
\end{exercise}

\begin{exercise}[$\star\star\star$]\label{exo:s2:storage-transport-and-physical-optionality:8}
\emph{Find the flaw.} ``Winter gas has been dearer than summer gas in most years, so we lease storage, inject all summer and sell in winter; hedging
only costs us money.''
\end{exercise}

\section{Problem: The Asset as Options}

\begin{problem}[{Weekend problem --- storage, transport and physical optionality}]\label{pb:s2:storage-transport-and-physical-optionality:1}
The chapter's synthetic lease, pipe and cargo and the Henry Hub record.

\medskip\noindent\textbf{Part I --- Storage.}
\begin{enumerate}
\item Define a \omterm{def:s2:storage-transport-and-physical-optionality:storage}{storage deal}.
\item Describe the lease, the curve and the intrinsic plan.
\item Give the four values and what each method assumes.
\item What did Secomandi and Lai, Margot and Secomandi find?
\end{enumerate}

\medskip\noindent\textbf{Part II --- Monetising.}
\begin{enumerate}[resume]
\item Describe the static and rolling programmes.
\item Why does the static programme earn the same on every path?
\item Give the rolling programme's results at each cost.
\item What does gating do?
\end{enumerate}

\medskip\noindent\textbf{Part III --- Transport and cargoes.}
\begin{enumerate}[resume]
\item Define a \omterm{def:s2:storage-transport-and-physical-optionality:transport}{transport option}.
\item Value the transport strip, intrinsically and as options.
\item How does its value depend on correlation?
\item Value the cargo switch.
\end{enumerate}

\medskip\noindent\textbf{Part IV --- The verdict.}
\begin{enumerate}[resume]
\item State the \emph{named result}: storage value by method and the extra value of rolling the hedge.
\item What did an unhedged operator earn at Henry Hub from 2007 to 2025?
\item Define \omterm{def:s2:storage-transport-and-physical-optionality:assetbacked}{asset-backed trading}; what does the trader own?
\item Why is least-squares Monte Carlo below rolling intrinsic here?
\item Why do the spread-option pairs add so little here?
\item What would you pay for the lease?
\item Which strategy file would fail first if the forward curve went flat?
\item In one sentence: what does a storage trader sell?
\end{enumerate}
\end{problem}

\section{Interview questions}

\begin{interviewq}[$\star$ \iqroles{trader}]\label{iq:s2:storage-transport-and-physical-optionality:1}
What is the intrinsic value of a storage facility, and what is extrinsic value?
\end{interviewq}

\begin{interviewq}[$\star\star$ \iqroles{researcher}]\label{iq:s2:storage-transport-and-physical-optionality:2}
How would you value gas storage, and how would you know your method is close to optimal?
\end{interviewq}

\begin{interviewq}[$\star\star$ \iqroles{trader}]\label{iq:s2:storage-transport-and-physical-optionality:3}
The winter-summer spread widens by 30 cents. What do you do with your storage hedges?
\end{interviewq}

\begin{interviewq}[$\star\star$ \iqroles{risk}]\label{iq:s2:storage-transport-and-physical-optionality:4}
How would you measure the risk of a storage book that is fully hedged at intrinsic?
\end{interviewq}

\begin{interviewq}[$\star\star$ \iqroles{developer}]\label{iq:s2:storage-transport-and-physical-optionality:5}
The rolling-intrinsic re-optimisation must run on every path every month. How would you make it fast?
\end{interviewq}

\begin{interviewq}[$\star\star\star$ \iqroles{researcher}]\label{iq:s2:storage-transport-and-physical-optionality:6}
Show that the static programme's P\&L equals the intrinsic value on every path, whatever the spot prices, and say which assumptions this needs.
\end{interviewq}
